\subsection{Cayley's Formula}

Cayley's Formula counts the number of distinct labeled trees on $n$ vertices.

\subsubsection{Statement}

The number of labeled trees with $n$ vertices is:

\begin{equation}
    n^{n-2}
\end{equation}

\subsubsection{Special Cases}

\begin{center}
\begin{tabular}{c|c|c}
$n$ & $n^{n-2}$ & Trees \\
\hline
1 & $1^{-1} = 1$ & 1 \\
2 & $2^0 = 1$ & 1 \\
3 & $3^1 = 3$ & 3 \\
4 & $4^2 = 16$ & 16 \\
5 & $5^3 = 125$ & 125 \\
6 & $6^4 = 1296$ & 1296 \\
\end{tabular}
\end{center}

\subsubsection{Prüfer Code}

Cayley's Formula can be proven using the Prüfer code, which establishes a bijection between:

\begin{itemize}
    \item Labeled trees with $n$ vertices
    \item Sequences of length $n-2$ with elements from $\{1, 2, \ldots, n\}$
\end{itemize}

There are exactly $n^{n-2}$ such sequences.

\subsubsection{Generalizations}

\textbf{Forests:} The number of labeled forests with $n$ vertices and $k$ trees is:

\begin{equation}
    k \cdot n^{n-k-1}
\end{equation}

\textbf{Spanning trees:} For a complete graph $K_n$, the number of spanning trees equals $n^{n-2}$ (which is Cayley's Formula).

\subsubsection{Applications}

Cayley's Formula is used in:

\begin{itemize}
    \item Counting problems involving labeled trees
    \item Graph enumeration
    \item Network design and reliability analysis
    \item Combinatorial optimization
\end{itemize}
